\documentclass{academics}

\academicsCourse{Java 与 Web 开发}{Java and Web Development}
\academicsReport{实验三}{2026 年 4 月 20 日}
\academicsArthur{华博文}{19240212}

\begin{document}

\makeacademicscover

\section{实验环境}

\subsection{操作系统}

macOS 26.4 arm64，Darwin 25.4.0，Apple M5。

\subsection{编程工具}

Visual Studio Code 1.117.0，NeoVim v0.12.0，openjdk 21.0.10。

\subsection{其他工具}

\LaTeX{}，\mil{mjr.py}。

\section{实验内容}

\subsection{第 1 题}

\paragraph{题目描述} {
    编写一个 Java 程序，实现一个简单的除法计算器。要求：使用 \mil{Scanner} 从控制台读取两个整数。计算第一个数除以第二个数的结果（整数除法），并输出。

    处理可能出现的异常：
    \begin{itemize}
        \item 如果用户输入的不是整数，捕获 \mil{InputMismatchException}，输出``输入错误：请输入整数！''。
        \item 如果除数为 $0$，捕获 \mil{ArithmeticException}，输出``数学错误：除数不能为零！''。
    \end{itemize}
    无论是否发生异常，最后都要输出``程序结束''。
}

\paragraph{程序实现} {
    \inputminted[label=DivisionCalculator.java]{java}{../src/P01/DivisionCalculator.java}
}

\paragraph{运行结果} {
    \inputminted{text}{../res/P01.out}
}

\subsection{第 2 题}

\paragraph{题目描述} {
    定义一个字符串数组 \mil[java]{String[] arr = { "10", "20", "abc", "30" }}。编写一个方法 \mil[java]{getIntegerFromArray(int index)}，接收一个索引，尝试将数组中对应位置的元素转换为整数并返回。方法需要处理：

    \begin{itemize}
        \item 如果索引超出数组范围，抛出 \mil{ArrayIndexOutOfBoundsException}，方法内捕获该异常并输出``索引越界，有效范围为0~3''。
        \item 如果字符串无法转换为整数，抛出 \mil{NumberFormatException}，方法内捕获并输出``数组元素不是有效整数''。
        \item 无论是否发生异常，最后都要输出``方法执行完毕''（在 \mil{finally} 中实现）。
    \end{itemize}
    在 \mil{main} 方法中循环调用该方法，传入索引 $0$、$1$、$2$、$3$、$4$，观察输出。
}

\paragraph{程序实现} {
    \inputminted[label=ArrayConverter.java]{java}{../src/P02/ArrayConverter.java}
}

\paragraph{运行结果} {
    \inputminted{text}{../res/P02.out}
}

\subsection{第 3 题}

\paragraph{题目描述} {
    定义一个自定义检查型异常类 \mil{InvalidAgeException}，继承 \mil{Exception}。
    编写一个 \mil{Person} 类，包含私有属性 \mil{age}，提供构造方法和 \mil[java]{setAge(int age)} 方法。在 \mil{setAge} 中，如果年龄小于 $0$ 或大于 $150$，则抛出 \mil{InvalidAgeException}，异常消息为``年龄必须在0~150之间''。

    在测试类的 \mil{main} 方法中，创建一个 \mil{Person} 对象，尝试设置一个非法年龄（如 $-5$），并捕获异常，打印异常信息。然后尝试设置合法年龄（如 $25$），并输出设置成功。
}

\paragraph{程序实现} {
    \inputminted[label=InvalidAgeException.java]{java}{../src/P03/InvalidAgeException.java}
    \inputminted[label=Person.java]{java}{../src/P03/Person.java}
    \inputminted[label=AgeTest.java]{java}{../src/P03/AgeTest.java}
}

\paragraph{运行结果} {
    \inputminted{text}{../res/P03.out}
}

\subsection{第 4 题}

\paragraph{题目描述} {
    定义一个自定义运行时异常类 \mil{InsufficientBalanceException}，继承 \mil{RuntimeException}。
    编写一个 \mil{BankAccount} 类，包含私有属性 \mil{balance}（余额）。提供构造方法初始化余额，以及 \mil[java]{withdraw(double amount)} 方法。在 \mil{withdraw} 中，如果取款金额大于余额，则抛出 \mil{InsufficientBalanceException}，异常消息为``余额不足！当前余额：xx，取款金额：xx''；否则扣除相应金额。

    在测试类的 \mil{main} 方法中，创建一个余额为 $1000$ 的账户，尝试取款 $1200$，捕获异常并打印信息。然后再取款 $500$，打印取款后的余额。
}

\paragraph{程序实现} {
    \inputminted[label=InsufficientBalanceException.java]{java}{../src/P04/InsufficientBalanceException.java}
    \inputminted[label=BankAccount.java]{java}{../src/P04/BankAccount.java}
    \inputminted[label=BankTest.java]{java}{../src/P04/BankTest.java}
}

\paragraph{运行结果} {
    \inputminted{text}{../res/P04.out}
}

\subsection{第 5 题}

\paragraph{题目描述} {
    定义异常类 \mil{DataAccessException}（检查型异常，继承 \mil{Exception}），用于表示数据访问层错误。
    编写一个工具类 \mil{FileUtil}，包含静态方法 \mil[java]{readFileContent(String fileName)} 声明抛出 \mil{IOException}（在方法内模拟抛出）。再编写一个业务类 \mil{DataService}，包含方法 \mil[java]{getDataFromFile(String fileName)}，它内部调用 \mil[java]{FileUtil.readFileContent}。在 \mil{getDataFromFile} 中捕获可能抛出的 \mil{IOException}，然后将其包装为 \mil{DataAccessException} 并重新抛出（异常链，保留原始原因）。

    在测试类中调用 \mil{getDataFromFile} 并捕获 \mil{DataAccessException}，打印异常信息和原始异常堆栈。可以使用 \mil[java]{throw new DataAccessException("读取文件失败", cause)} 构造异常链。为了模拟，在 \mil{FileUtil.readFileContent} 中直接 \mil[java]{throw new IOException("文件不存在")}。
}

\paragraph{程序实现} {
    \inputminted[label=DataAccessException.java]{java}{../src/P05/DataAccessException.java}
    \inputminted[label=FileUtil.java]{java}{../src/P05/FileUtil.java}
    \inputminted[label=DataService.java]{java}{../src/P05/DataService.java}
    \inputminted[label=DataServiceTest.java]{java}{../src/P05/DataServiceTest.java}
}

\paragraph{运行结果} {
    \inputminted{text}{../res/P05.out}
}

\subsection{第 6 题}

\paragraph{题目描述} {
    编写一个父类 \mil{Parent}，其中定义方法 \mil[java]{void readFile() throws IOException}（方法体可为空或简单输出）。
    编写两个子类：\mil{Child1} 重写 \mil{readFile} 方法，声明抛出 \mil{FileNotFoundException}（\mil{IOException} 的子类）；\mil{Child2} 重写 \mil{readFile} 方法，不声明抛出任何异常。

    分别创建子类对象，并通过父类引用调用 \mil{readFile} 方法，观察编译和运行情况。在 \mil{main} 中，使用 \mil{try-catch} 处理可能抛出的异常。

    要求：解释为什么这样声明是合法的（或非法），并编写代码验证。子类方法可以抛出更具体的异常或不抛出异常，但不能抛出父类方法未声明的更宽泛的异常（如 \mil{Exception}）。
}

\paragraph{程序实现} {
    \inputminted[label=Parent.java]{java}{../src/P06/Parent.java}
    \inputminted[label=Child1.java]{java}{../src/P06/Child1.java}
    \inputminted[label=Child2.java]{java}{../src/P06/Child2.java}
    \inputminted[label=OverrideTest.java]{java}{../src/P06/OverrideTest.java}
}

\paragraph{运行结果} {
    \inputminted{text}{../res/P06.out}
}

\subsection{第 7 题}

\paragraph{题目描述} {
    编写一个类 \mil{FinallyReturnDemo}，包含一个静态方法 \mil[java]{int calculate()}，其中：
    \begin{itemize}
        \item 在 \mil{try} 块中，打印``进入try''，返回 $1$。
        \item 在 \mil{catch} 块中（可包含一个永远不会发生的异常），打印``进入catch''，返回 $2$。
        \item 在 \mil{finally} 块中，打印``进入finally''，返回 $3$。
    \end{itemize}
    在 \mil{main} 方法中调用 \mil{calculate()} 并打印返回值，观察输出结果，解释为什么。修改：将 \mil{finally} 中的 \mil{return} 注释掉，再观察返回值。

    理解 \mil{finally} 块中的 \mil{return} 会覆盖 \mil{try} 或 \mil{catch} 中的 \mil{return}，以及 \mil{finally} 总是在 \mil{return} 之前执行。
}

\paragraph{程序实现} {
    \inputminted[label=FinallyReturnDemo.java]{java}{../src/P07/FinallyReturnDemo.java}
}

\paragraph{运行结果} {
    \inputminted{text}{../res/P07.out}
}

\subsection{第 8 题}

\paragraph{题目描述} {
    定义一个自定义检查型异常 \mil{BusinessException}，继承 \mil{Exception}。该类包含：
    \begin{itemize}
        \item 私有属性 \mil{errorCode}（\mil[java]{int}），表示错误码。
        \item 构造方法 \mil[java]{BusinessException(String message, int errorCode)}，调用父类构造器，并设置错误码。
        \item 提供 \mil[java]{getErrorCode()} 方法。
    \end{itemize}
    编写一个业务类 \mil{OrderService}，包含方法 \mil[java]{void placeOrder(int quantity)}，如果 \mil{quantity} 小于等于 $0$，则抛出 \mil{BusinessException}，错误码为 $1001$，消息为``订单数量必须大于0''；如果 \mil{quantity} 大于 $100$，则抛出 \mil{BusinessException}，错误码为 $1002$，消息为``订单数量不能超过100''。

    在测试类中调用 \mil{placeOrder}，捕获 \mil{BusinessException}，输出错误码和异常消息。
}

\paragraph{程序实现} {
    \inputminted[label=BusinessException.java]{java}{../src/P08/BusinessException.java}
    \inputminted[label=OrderService.java]{java}{../src/P08/OrderService.java}
    \inputminted[label=OrderTest.java]{java}{../src/P08/OrderTest.java}
}

\paragraph{运行结果} {
    \inputminted{text}{../res/P08.out}
}

\subsection{第 9 题}

\paragraph{题目描述} {
    定义一个 \mil{Person} 类，包含私有属性 \mil{name}（\mil{String}）和 \mil{age}（\mil[java]{int}）。编写构造方法 \mil[java]{Person(String name, int age)}，要求：
    \begin{itemize}
        \item 如果 \mil{name} 为 \mil{null} 或空字符串（长度为 $0$），抛出 \mil{IllegalArgumentException}，异常消息为``姓名不能为空''。
        \item 如果 \mil{age} 小于 $0$ 或大于 $150$，抛出 \mil{IllegalArgumentException}，异常消息为``年龄无效：$x$''（$x$ 为实际年龄）。
    \end{itemize}
    注意：\mil{IllegalArgumentException} 是运行时异常，无需显式声明。

    在测试类的 \mil{main} 方法中，分别尝试创建合法对象和非法对象，捕获异常并打印异常信息。对于非法情况，观察程序是否继续执行。
}

\paragraph{程序实现} {
    \inputminted[label=Person.java]{java}{../src/P09/Person.java}
    \inputminted[label=ConstructorTest.java]{java}{../src/P09/ConstructorTest.java}
}

\paragraph{运行结果} {
    \inputminted{text}{../res/P09.out}
}

\end{document}
